\subsection*{14.4 Central Limit Theorems}

We use Levy's continuity theorem to prove a basic version of central limit theorem,

which assumes iid. random variables and finite variance.

\begin{thmbox}{14.7 (Lindeberg-Levy Central Limit Theorem)}
\end{thmbox}

Let .(Xk)\infty

k=1 be a sequence of iid. random variables with finite mean

E[X1]= \mu and finite variance Var(X1)= \sigma 2. Then

\sumn

k=1 Xk -n\mu

\sigma\sqrt n converges in distribution toN( 0,1).

Sketch of Proof We may assume \mu= 0 by re-defining X1 :=X 1 -\mu Let. Sn \coloneqq

X1 +X 2 +\cdot\cdot\cdot+ Xn.

Step 1. [ 4, Theorem 3.3.20] Since. X1 has variance. \sigma2, the characteristic function

of X1 is

\phiX1 (t)= 1 - \sigma2t2

2 +o(t 2),

where o(t2) refers to a function h(t) such that limt\to 0 h(t)

t2 =0 .

Step 2. Because the random variables Xk's are assumed to be independent, the

characteristic function of Sn

\sigma\sqrt n is

(

1- t2

2n +o

( 1

n

)) n

It can be written as .

(

1+ cn

n

) n

, where cn be the constant defined by

cn

n =- t2

2n +o

( 1

n

)

Step 3. [ 4, Theorem 3.4.2] If c1, c2, c3,.are complex numbers with

limn\to\infty cn =c , then

limn\to\infty (1+ c n/n)n =e c.

Step 4. Because cn \to- t2

2 +o( 1), we get limn\to\infty cn =- t2/2. The charac-

teristic function of Sn

\sigma\sqrt n approaches e-t2/2 as n\to\infty Since e-t2/2 is a continuous

function, by Levy's continuity theorem, Sn

\sigma\sqrt n converges to the distribution whose

characteristic function is e-t2/2.

Step 5. Steps 1-4 hold for any iid. sequence with finite mean and variance.

In particular, it is true if the random variables are iid. standard Gaussian random

variables. We use the fact that the sum of Gaussian random variables are Gaussian.

When Steps 1 and 4 are applied to iid. standard Gaussian random variables, the

random variables Sn

\sigma\sqrt n have standard Gaussian distribution for all n and hence must

converge in distribution to N( 0,1)The characteristic function of N( 0,1) must be

equal to e-t2/2. By the uniqueness of characteristic function, Sn

\sigma\sqrt n converges in

distribution to N( 0,1)\hfill $\square$

Based on the central limit theorem, we have the following approximations of

binomial and Poisson distributions using normal distribution.

\textbf{Example 14.4.1 (Normal Approximation of Binomial Distribution)} \\

Suppose Xn's are iid. Bernoulli random variables with success probability p The mean and

variance of. Xn are p andp(1- p), respectively. The sumSn =X 1 +X2 +\cdot\cdot\cdot+ Xn has distribution

Binom(n, p)By the central limit theorem,

Sn -np\sqrt np(1- p)

D

-\to N( 0,1).

\textbf{Example 14.4.2 (Normal Approximation of Poisson Distribution)} \\

We use the fact that the sum of n iid. Poisson random variables with mean 1 is a Poisson random

variable with mean nLet. Sn denote a Poisson random variable with mean n. Since the variance

of a Poisson random variable with mean 1 is 1, we can apply the central limit theorem to get

(Sn -n)/ \sqrt n

D

-\to N( 0,1).

While the Lindeberg-Levy central limit theorem is a useful result, it is limited

in that it requires the random variables to be iid. However, it is possible to relax

the iid. assumption and still obtain similar results. In the followings, we will state

a central limit theorem for triangular array.

For each positive integer n ,l e t kn be a positive integer and let Xn1,

Xn2,...,X n,kn be independent random variables defined on a common probability

space. We can represent these random variables as a collection of rows, where the

n-th row contains kn random variables.

X11,X 12,...,X 1,k1

X21,X 22,X 23,...,X 2,k2

X31,X 32,X 33,X 34,...,X 3,k3

..

Note that the probability spaces may be different for different rows.

Suppose that the expectations of all random variables are zero, for notational

simplicity, and let the variance of Xn,i be \sigma2

ni ,f o rn\geq 1 and .1\leq i \leq n W el e t

Sn be the sum of the kn random variables in the n-th row, and s2

n = \sumkn

i=1 \sigma2

ni ,f o r

n\geq 1 , be the variance of Sn. We assume that s2

n \to\infty as n\to\infty .

\begin{thmbox}{14.8 (Central Limit Theorem for Triangular Array)}
\end{thmbox}

With the notation above, we have

Xn1 +X n2 +\cdot\cdot\cdot+ Xn,kn

sn

D

-\to N( 0,1)

if one of the following conditions hold:

(Lindeberg condition) For all\epsilon> 0,

s2n

kn\sum

i=1

E[X2

n,i 1{\vertXn,i\vert\geq\epsilonsn}}]\to 0as n \to\infty .

(Lyapunov condition) There exists\delta> 0 such that

s2+\deltan

n\sum

i=1

E[\vertXn,i\vert2+\delta]\to 0as n \to\infty .

In the Lindeberg condition, the expectation is the integral

\int

{\vertx\vert\geq\epsilonsn}

x2 d\mun,i(x),

integrating the function x2 in the range \vertx\vert\geq \epsilonsn with respect to the measure

\mun,i induced by Xn,i Intuitively, this means that the tail of each Xk must decrease

fast enough so that the sum of the tail parts of the variances does not contribute

significantly to the overall variance of Sn. We remark that for independent random

variables, the conditions in the Lindeberg-Levy central limit theorem imply the

Lindeberg condition.

On the other hand, the Lyapunov condition is a stronger condition than the

Lindeberg condition. Specifically, the Lyapunov condition requires that the tail

probability of each random variable Xk decays sufficient fast so that certain power

law holds. This condition is useful in some situations where the Lindeberg condition

is difficult to check.

The proof is beyond the scope of this book. We refer the readers to [ 2] and [4].

Instead of going through the proof, we give an example and demonstrate how to

apply this theorem.

\textbf{Example 14.4.3 (Asymptotic Normality in Coupon Collection Problem)} \\

In the coupon collection problem, we consider the random experiment that repeatedly and

independently draws one out of n different coupons, where n is a fixed parameter. LetTn,m denote

the first time that exactly m distinct coupons have been collected. This random variable is the sum

of m independent geometric-distributed random variables,

Tn,m =X n,0 +X n,1 +X n,2 +\cdot\cdot\cdot+ Xn,m-1.

The success probability of random variableXn,i is. 1- i

n ,a n dXn,0 is equal to 1 with probability 1.

For i\geq 1 ,Xn,i counts the number of additional coupons needed to be drawn in order to get a new

one after i distinct coupons have been collected.

The pmf of the random variableXn,i is given by

P(Xn,i =j) = ( 1- p n,i)j- 1pn,i

for j\geq 1 ,w h e r epn,i \coloneqq1 - i/n is the probability of success. By using the moment generating

function,

E[eXn,it ]= pn,iet

1- ( 1- p n,i)et ,

we can compute the mean and variance ofXn,i as

E[Xn,i]= 1

pni

,Var (X 2

n,i)= 1- p n,i

p2

n,i

Additionally, the centered fourth moment ofXn,i is given by

E[(Xn,i -E [Xn,i])4]= 1

p4

n,i

(1- p n,i)(p2

n,i -9 pn,i +9 ).

Suppose we increase n and m simultaneously such that the ratio m/n approaches 1/2. For

instance, we may take m= (n+ 1 )/2. We are interested in the asymptotic distribution of the

random variableTn,(n+1)/2.

We can write the random variables Xn,i ,f o rn\geq 2 and .0\leq i \leq (n+ 1 )/2, as a triangular

array. The random variables in each row are independent, but not identically distributed. Hence,

we cannot apply Theorem 14.7. Instead, we will apply Theorem 14.8, with\delta= 2 in the Lyapunov

condition. We obtain the following upper bound for the sum of fourth moments,

s4n

(n+1)/2\sum

i=0

E[(Xn,i -E [Xn,i])4])]= 1

s4n

(n+1)/2\sum

i=0

p4

n,i

(1- p n,i)(p2

n,i -9 pn,i +9 )

\leq 1

s4n

(n+1)/2\sum

i=0

p4

n,i

= 9

s4n

(n+1)/2\sum

i=0

(1- i/n) 4 .

We approximate the sum by a Riemann integral,

s4n

(n+1)/2\sum

i=0

E[(Xn,i -E [Xn,i])4])]\leq 9n

s4n

\int 1/2

(1- x) 4 dx. (14.9)

Because

s2

n =

(n+1)/2\sum

i=0

i/n

(1- i/n) 2 n

\int 1/2

(1- x) 2 dx= n( 1- log (2)),

we see that right-hand side of (14.9) approaches 0 as n approaches infinity. This verifies the

Lyapunov condition. By using approximations by Riemann integral, we can show that the mean of

E[Tn,(n+1)/2]\to nlog 2. We apply Theorem 14.8 to conclude that, as n approaches infinity,

Tn,(n+1)/2-n log 2\sqrt

n(1- log 2)

D

-\to N( 0,1).

We implement the coupon collection problem using the following Python program and

generate the histogram of the number of coupons we draw until we obtainn/2 distinct coupons out

of n coupons.

from random import randint

from numpy import log, sqrt

import matplotlibpyplot as plt

def coupon_collection(n,m):

L = [] # intitialize to empty list

while True:

Lappend(randint(1,n)) # generate a new coupon

if len(set(L)) >= m:

break # break if we have m distinct coupons

return(len(L))

n = 100 # there are n types of coupons

m = 50 # stop when m distinct coupons are collected

K = 50000 # repeat the experiment K times

T = [coupon_collection(n,m) for j in range(K)] # generate data

m u=n *log(2)

sigma2 = n *(1-log(2))

x_min , x_max = (50,100) # range of plots

bins = [i for i in range(x_min,x_max)] # histogram bins

plthist(T,bins=bins) # plot histogram

X = [x for x in range(x_min,x_max)]

Y = [K/sqrt(2 *pi*sigma2)*exp(-(x-mu)**2/(2*sigma2)) for x in X]

pltplot(X,Y, 'k-')

pltxlabel('Number of coupons drawn until we get 50 out of 100 coupons')

We plot the histograms for n= 100 and n= 1000 The figures also show the probability

density function of the limiting Gaussian distribution as a reference. The simulation results confirm

that the probability distribution of the number of draws until obtainingn/2 distinct coupons is well

approximated by a Gaussian distribution.
